\documentclass[10pt]{article}
\usepackage[margin=20mm]{geometry}
\usepackage{amsmath,amssymb,mathtools,fontspec,unicode-math,xurl,hyperref}
\IfFontExistsTF{DejaVu Serif}{\setmainfont{DejaVu Serif}}{\setmainfont{Latin Modern Roman}}
\setmathfont{latinmodern-math.otf}
\setlength{\parindent}{0pt}\setlength{\parskip}{6pt}\sloppy
\title{Technical verification of
\(E_\infty\)-cells and general linear groups}
\author{AI-assisted technical review}
\date{}
\begin{document}
\maketitle
Abstract: The acquired paper is examined through its building models, invariant Steinberg pairing, bar constructions and stability applications. A rational three-dimensional example refutes the local product identification used in the semi-local connectivity proof. Two further core derivations require correction: the partial-order group comparison and the grading of the \(3\)-local argument for Steinberg homology. A further quantitative rank-filtration quotient is false. Local findings concern the ring pre-Bloch presentation and two Steinberg-module boundary assertions. The field-case divided-power calculation and the stated spectral-sequence ranges are supported within the inspected dependencies. Major is the supported overall assessment. This is an initial technical review, with scientific adjudication pending.\par
\section*{Identity and scope}
This report concerns Søren Galatius, Alexander Kupers and Oscar Randal-Williams, \(E_\infty\)-cells and general linear groups of infinite fields, arXiv:2005.05620v3 [GKRW]. The primary PDF has 82 pages, an arXiv watermark dated 18 January 2025 and a printed date of 22 January 2025. The frozen journal identity is Duke Mathematical Journal 174 (2025), 2927–3046, DOI https://doi.org/10.1215/00127094-2025-0005. A separate Cambridge accepted manuscript [CAM] has 101 pages and is dated 17 January 2025. Neither acquired file is certified identical to the 120-page journal version.\par
The complete 5,027-line primary PDF extraction was read, with targeted checks against the actual PDF pages and the released source. The review covers Sections 1–10, the hypotheses and dependencies of Theorems A–H, the finite-rank and coefficient restrictions, and the figures describing low-degree homology. It does not reprove the full external theory of cellular \(E_k\)-algebras or every cited theorem. Native bounded checks were integrated into the primary analysis. All page references without a version prefix refer to one-based primary PDF pages; its printed numbering agrees. The accepted manuscript is cited by its own PDF pages.\par
The primary convention is that rings are commutative and that the semi-local applications require connectedness and infinite residue fields. The rank-zero part of the non-unital algebra is empty. A statement proved for infinite fields is not silently extended to finite fields. Group homology and algebraic \(K\)-theory use discrete groups unless a topological-ring comparison is expressly invoked.\par
\section*{Supported parts of the argument}
The Steinberg tensor-square coinvariant calculation in Section 2 uses the Lee–Szczarba relations and successive Jordan-block reductions. The pairing of a standard apartment with the one-block apartment is \(1\), so cyclicity together with the pairing gives the asserted integral isomorphism. In Section 6.3 the sign \(\epsilon(n)=n(n-1)/2\) obeys \[\epsilon(m+n)=\epsilon(m)+\epsilon(n)+mn.\] Consequently the apartment-product calculation gives \[\gamma_m(x)\gamma_n(x)=\binom{m+n}{m}\gamma_{m+n}(x).\] The factorial cancellation takes place over the torsion-free integral coinvariants, before base change; it does not divide by a factorial in an arbitrary coefficient ring. The prime-power boundary groups in Theorem B then agree with the indecomposables of this divided-power algebra. The case \(n=1\) is the free generator, handled separately.\par
For a field, two flags give a full two-dimensional lattice by taking intersections, so the global identification \(D^2(V)=D^1(V)\wedge D^1(V)\) in Proposition 6.3 is supported. Problem 1 concerns different, restricted subspaces indexed by a fixed simplex \(\sigma\); it does not contradict that global field statement.\par
Conditional on the asserted \(E_2\) vanishing, the passage to \(E_\infty\) vanishing preserves the line \(d<2n-2\). Indeed, in a further reduced bar construction of length \(p\), positive ranks \(n_1+\cdots+n_p=n\) contribute degree at least \(2n+p(k-1)\) before desuspending by \(k+1\). The resulting lower bound is \[2n-2+(p-1)(k-1)\geq 2n-2\qquad(k\geq2).\] For fixed \(n\), only \(1\leq p\leq n\) occurs. Thus the formal bar argument is not the source of the connectivity defect below.\par
The Harrison comparison in Section 9.3 has the announced thresholds \(n\geq5\) for semi-local rings and \(n\geq4\) for fields. The rational and mod-\(p\) minimal-degree arguments preserve their strict slopes \((4n-1)/3\), \(3n/2\) and \(5n/3\) under the stated inputs. The factorial range in Section 9.6 was checked separately from the exception \(n=3=\operatorname{char}(k)\); the paper explicitly imports the known low-rank result for that exception. These checks are conditional on the preceding inputs and do not erase the defects or the unresolved dependency recorded below.\par
\section*{Confirmed problems}
Equation reference: eq:local-product. Problem 1 (A local product isomorphism fails for compatible lattices).  Proposition 7.24, pp. 43–44, asserts for every simplex \(\sigma<E(M)\) a bijection \[
 D^2(M)^\sigma_{p,q}\longrightarrow
 D^1(M)^\sigma_p\wedge D^1(M)^\sigma_q.\] Definitions 7.14–7.15, pp. 38–39, require more than membership of each lattice entry in \(\operatorname{Sub}(M)_\sigma\): all the sums in equation (7.6) must belong to that set. Lemma 7.16 makes this distinction explicit.\par
Take the connected semi-local ring \(A=\mathbb Q\), the free module \(M=\mathbb Q^3\) with basis \(e_1,e_2,e_3\) and dual basis \(x,y,z\), and \[\sigma=\{x,y,z,x+y+z\}.\] Every member is a surjection to \(\mathbb Q\), and every subset has projective cokernel because it is a map of vector spaces. Thus \(\sigma\) is an actual simplex of \(E(M)\) under Definition 7.5, with no additional genericity assumption required. Put \[P=\ker x\cap\ker y=\mathbb Qe_3,\qquad
 Q=\ker z\cap\ker(x+y+z)=\mathbb Q(e_1-e_2).\] The flags \(0<P<M\) and \(0<Q<M\) both lie in \(D^1(M)^\sigma_2\). The only two-dimensional lattice with these boundary flags has \[\phi(1,1)=0,\quad \phi(1,2)=P,\quad
 \phi(2,1)=Q,\quad \phi(2,2)=M,\] and is zero when either coordinate is zero. It is a full lattice: \(P\cap Q=0\), and \(P+Q\) is a direct summand. For example, a decomposition as in Lemma 5.6 is obtained by putting \(P\) and \(Q\) in positions \((1,2)\) and \((2,1)\) and choosing a one-dimensional complement in position \((2,2)\).\par
At \(a=(1,1)\), however, equation (7.6) requires \(P+Q\in\operatorname{Sub}(M)_\sigma\). Its annihilator is \[(P+Q)^\perp=\mathbb Q(x+y).\] If \(P+Q\) were an intersection of kernels of a subset of \(\sigma\), that subset would span \(\mathbb Q(x+y)\). None of the four elements of \(\sigma\) belongs to this line. The empty subset spans zero, and no nonempty subset can span that line. This also gives an exact finite check over all 16 subsets. Thus the displayed pair of flags has no preimage in \(D^2(M)^\sigma_{2,2}\), disproving [eq:local-product] (\#eq:local-product).\par
The proof on p. 43 verifies the lattice condition, but does not verify this stronger compatibility requirement. Lemma 7.10 only supplies projectivity of a relevant quotient; it does not put the sum back in \(\operatorname{Sub}(M)_\sigma\). The warning immediately before Lemma 7.16 that this set need not be closed under sums is consistent with the example.\par
Impact and repair boundary. The proof of Theorem 7.1(ii) on p. 44 starts with this claimed isomorphism and deduces \((2d-1)\)-connectivity. The proof of Theorem 7.6(7.2) on the same page also uses it. These feed the semi-local \(E_2\) calculation, Corollary 7.3 and Theorem A, and hence the semi-local applications. This is a Major invalid derivation closely tied to the main contribution. It is not a counterexample to Theorem A or to global field connectivity. Appendix A verifies a repair by changing the restricted pieces to require only entrywise membership and proving a new homotopy-colimit comparison. It restores this connectivity and resolution step under the original semi-local hypotheses. Lemmas 7.16–7.18 are not carried over to the changed pieces. This is a replacement argument; the original Proposition 7.24 remains false. The same original definitions and failed product claim occur on accepted pp. 47–49 and 53–54.\par
Problem 2 (The \(3\)-local proof uses incompatible degrees).  In the proof of Theorem 10.2, pp. 77–78, the model in (10.2) uses \(b\) of bidegree \((2,4)\) and a cell \(\rho_3\) of bidegree \((3,5)\) attached along \((\beta Q^1_3)_\mathbb Z(\sigma)-\sigma b\). After one bar construction, (10.3) and the next displays place \(\bar b\) in \((2,5)\) and the attached \(\bar\rho_3\) cell in \((3,6)\). The subsequent spectral sequence nevertheless treats \(\bar\rho_3\) as an exterior generator.\par
The earlier operation calculation, Lemma 9.2(ii) on p. 56, puts \(b\) in \(H_{2,3}(\mathbf R)\) and \((\beta Q^1_3)_\mathbb Z(\sigma)\) in \(H_{3,3}(\mathbf R)\). In the printed universal model, the two proposed attaching terms have different homological degrees, \(3\) and \(4\). They therefore do not define the stated homogeneous attaching map. Moreover a generator of the printed degree \(6\) is even, whereas the displayed Koszul pair uses it as odd. These discrepancies affect the actual differential and the cancellation of the low-slope polynomial generator; they are not changes to a citation or to a name alone.\par
A coherent correction has the following degrees: \[\begin{array}{c|ccccc}
 \text{object}&b&\rho_3\text{ cell}&\bar b&
       \beta Q^1_3(\bar\sigma)&\bar\rho_3\text{ cell}\\ \hline
 (n,d)&(2,3)&(3,4)&(2,4)&(3,4)&(3,5).
\end{array}\] The homology suspension raises degree \(3\) to \(4\). Writing \(u=\beta Q^1_3(\bar\sigma)\) and \(v=\bar\rho_3\), the relevant corrected differential graded factor over \(\mathbb F_3\) is \[\bigl(\Lambda_{\mathbb F_3}[\bar\sigma,v]\otimes\mathbb F_3[u],
        \ d v=u\bigr),\qquad |u|=(3,4),\quad |v|=(3,5).\] The factor \(\Lambda[v]\otimes\mathbb F_3[u]\) has homology \(\mathbb F_3\): multiplication by \(u\) is injective, and its cokernel is the constants. The remaining generators have slope at least \(3/2\) under the Dyer–Lashof degree estimates used in the source. Adding \(\bar b\) of degree \((2,4)\) preserves that bound. The motivating inequality also needs reversal: \((4k-1)/(3k-1)<3/2\) for \(k\geq2\), with equality for \(k=1\).\par
Impact and counterevidence. This supplies a corrected local calculation, conditional on the source’s cell-existence and connectivity inputs. The original degree-\(6\) exterior calculation is not that calculation. The affected argument proves the integral, including \(3\)-primary, range in Theorem 10.2, equivalently the named Theorem H. It is therefore Major as a demonstrated invalid core derivation, even though a coherent degree repair is available. No counterexample to the homology-vanishing conclusion is claimed. Both acquired PDF versions contain the incompatible degrees (accepted manuscript pp. 94–96).\par
Problem 3 (The partial-order comparison uses the wrong invariant module).  Proposition 5.13, pp. 21–22, proves a group-homology comparison by removing a maximal element \(t\) of a finite poset. The proof asserts \(f(M_t)\subset M_t\) and identifies a kernel with the full ambient group \(\operatorname{GL}(M,\operatorname{fix}M')\), where \(M'=\bigoplus_{x\ne t}M_x\).\par
Take \(A=\mathbb Q\), \(X=\{a,b,t\}\) with the only strict relation \(a<t\), and \(M_a=M_b=M_t=\mathbb Q\). These satisfy the Nesterenko–Suslin hypothesis. Under the paper’s convention that the matrix component \(M_x\to M_y\) is allowed when \(y\leq x\), the automorphism \(f=I+E_{a,t}\) belongs to \(\operatorname{GL}(M,\leq)\) but sends \(e_t\) to \(e_t+e_a\). Thus \(M_t\) is not invariant and the claimed induced action on \(M/M_t\) is not defined. In contrast, \(h=I+E_{b,t}\) fixes \(M'\) pointwise, so it belongs to the displayed ambient kernel group, but it does not belong to \(\operatorname{GL}(M,\leq)\) at all. The written sequence is not a sequence of the claimed subgroups.\par
Here is a repair with exactly the original hypothesis. The invariant module is \(M'\), and set \(U=\bigoplus_{x<t}M_x\), a direct summand preserved by \(G'=\operatorname{GL}(M',\leq)\). Restriction gives a split extension with kernel \[K=\operatorname{GL}(M_t)\ltimes\operatorname{Hom}_A(M_t,U).\] The associated-graded map is \[\operatorname{GL}(M,\leq)\longrightarrow G'\times\operatorname{GL}(M_t).\] Its map on kernels is the Nesterenko–Suslin comparison for \((P,Q)=(U,M_t)\). It is a homology isomorphism, equivariantly for \(G'\), whose action on \(\operatorname{GL}(M_t)\) is trivial. The comparison of extension spectral sequences gives the required homology isomorphism. Induction over the removed relations then proves Proposition 5.13.\par
The correction is verified at the group and extension level; it does not require that \(t\) be above every other element. Nevertheless the original derivation uses an undefined quotient action and an incorrect kernel. Theorem 5.18 directly invokes it to identify split and ordinary buildings, a central input to both the field and semi-local calculations. The affected derivation is Major under the core-dependency criterion; the repaired proposition is supported, and no false homology-equivalence theorem is asserted.\par
Problem 4 (A rank-filtration quotient is not the whole first layer).  Remark 8.4, p. 52 (accepted p. 64), identifies \(K_{2r}(A)/F_rK_{2r}(A)\) with the first quotient layer’s homotopy group, and then with \(H^{E_\infty}_{r+1,2r}(\mathbf R_\mathbb Z)\).\par
Choose the discrete field \(A=\mathbb C\) and \(r=1\). It is connected and semi-local with an infinite residue field, so it satisfies every ring hypothesis here. Corollary 6.12 gives \(H^{E_\infty}_{2,2}(\mathbf R_\mathbb Z)=\mathbb Z/2\). On the other hand Matsumoto’s theorem identifies \(K_2(\mathbb C)\) with Milnor \(K_2(\mathbb C)\); the paper itself recalls this on p. 64. For every integer \(m>0\), choose an \(m\)th root \(a^{1/m}\in\mathbb C\). Bilinearity of symbols gives \[\{a,b\}=m\{a^{1/m},b\}.\] Hence \(K_2(\mathbb C)\) is divisible. Every quotient of a divisible abelian group is divisible, while \(\mathbb Z/2\) is not. The stated equality is false.\par
Let \(C_r\) denote the cofiber of \(F_rK(A)\to K(A)\). The long exact sequence gives instead \[K_{2r}(A)/F_rK_{2r}(A)
 \ \cong\
 \ker\bigl(\pi_{2r}(C_r)\longrightarrow\pi_{2r-1}(F_rK(A))\bigr).\] The higher layers do not change \(\pi_{2r}(C_r)\) in the range under discussion, so this is a subgroup of the critical group, not necessarily the entire group. For \(\mathbb C,r=1\), the subgroup must be zero because it is both divisible and a subgroup of \(\mathbb Z/2\).\par
The weaker finiteness conclusion survives this correction. The adjacent Remark 8.5 expressly changes to topological rings and does not alter the discrete-field counterexample. The exact integral quotient is a quantitative application of the paper’s critical-group calculation to its central stable-rank-filtration question. The dependency is Theorem B, then the layer/Hurewicz identification, then the invalid replacement of a kernel by the whole layer. The claimed value changes from \(\mathbb Z/2\) to zero in the admissible example. This is Major for that application, even though Theorems A–B and the weaker finiteness consequence are not refuted.\par
Problem 5 (The ring pre-Bloch presentation admits inadmissible symbols).  Section 9.1.1, p. 53, defines \(\mathfrak p(A)\) with generators \(x\in A^\times\setminus\{1\}\). Equation (9.3), p. 55, then assigns a term \(x\wedge(1-x)\) in an exterior construction on \(A^\times\).\par
Let \(A=\mathbb Q[\varepsilon]/(\varepsilon^2)\) and \(x=1+\varepsilon\). The ring is connected and local with infinite residue field \(\mathbb Q\). The element \(x\) is a unit, with inverse \(1-\varepsilon\), and \(x\ne1\). Thus it is an allowed generator in the printed presentation, but \(1-x=-\varepsilon\) is not a unit. The displayed boundary expression is undefined.\par
This is not remedied by the five-term relations. When their stated parameters \(u,v,1-u,1-v,u-v\) are units, every argument appearing in the five terms has both itself and its complement invertible. Consequently the other generators never occur in a relation and form an additional free direct summand in the printed presentation. The required presentation uses \[W_A=\{x\in A: x\in A^\times,\ 1-x\in A^\times\}\] as its generator set. Over a field this coincides with the printed set.\par
The chain calculation with transverse rank-one summands already uses the admissible set, and the Milnor relations are correctly restricted to \(x,1-x\in A^\times\) in Section 9.2, p. 58. This is therefore a Minor mismatch in the stated ring presentation and boundary domain, with explicit consequences for the notation in Figure 1 and Theorem 9.1. It does not provide a counterexample to the standard admissible-symbol calculation or to the field stability theorems. The corrected domain must also be carried through the later presentations on pp. 61–62; simply applying the printed map to the additional generators is not possible. The same generator set is visible in the accepted version, p. 65.\par
Problem 6 (The finite-field extension of the adjoint argument is unsupported).  The paragraph introducing the proof of Corollary D and the proof itself, p. 11, use injectivity of the Steinberg pairing’s adjoint for a vector space over a field without an infinitude qualification. Theorem 2.4(iv), p. 10, only proves that injectivity for infinite fields, and Remark 2.5 explicitly warns against removing the assumption.\par
For the stated finite-field test take \(F=\mathbb F_2\), \(V=F^2\) and \(k=\mathbb F_3\). The coefficient ring is connected. There are three lines, so \[\operatorname{St}_k(V)=\ker(k^3\xrightarrow{\text{sum}}k).\] In the basis \(u=(1,-1,0)\), \(v=(0,1,-1)\), the restricted pairing has Gram matrix \[\begin{pmatrix}2&2\\2&2\end{pmatrix}.\] Its nonzero radical is spanned by \(u-v\). The \(k\)-linear endomorphism \(T(u)=u-v,\ T(v)=0\) is nonzero but is sent to the zero bilinear form. Thus the full \(\operatorname{Hom}_k\) injection used in the proof is false for these parameters.\par
This does not refute Corollary D. In this same example, the matrices of the two generating transpositions are \[\begin{pmatrix}-1&1\\0&1\end{pmatrix},\qquad
 \begin{pmatrix}1&0\\1&-1\end{pmatrix}.\] Their joint centralizer over \(\mathbb F_3\) consists of scalars. Thus the module is indecomposable and the equivariant endomorphism map remains injective here. The test map \(T\) is not equivariant. For infinite \(F\) the actual adjoint proof is valid, with the reference to part (iv). The finite-field extension needs another argument. This is a Minor scope defect outside the main infinite-field calculation; a false Corollary D is not established.\par
Problem 7 (The finite-dimensional-submodule assertion omits rank one).  The paragraph following Corollary D, p. 11, says that for a coefficient field with no nontrivial zero sum of squares, \(\operatorname{St}_k(V)\) has no finite-dimensional invariant submodules.\par
Take \(F=\mathbb Q\), \(V=\mathbb Q\) and \(k=\mathbb R\). The source explicitly allows nonzero finite-dimensional \(V\), including dimension one. Its reduced-homology convention gives \(\operatorname{St}_\mathbb R(\mathbb Q)=\mathbb R\); \(\operatorname{GL}_1(\mathbb Q)\) acts trivially. The whole module is therefore a nonzero one-dimensional invariant submodule. Even the infinite-base-field hypothesis is satisfied.\par
The displayed argument proves that such a finite-dimensional submodule is a direct summand. Indecomposability then makes it the whole module; it does not rule it out unless the whole Steinberg module is known to be infinite-dimensional. For infinite \(F\) and \(\dim V\geq2\), the apartment basis indexed by the upper-unipotent group supplies that missing fact. Adding these restrictions restores the intended argument. This local boundary assertion is Minor and does not affect Theorem C or the divided-power calculation.\par
\section*{Clarifications and unresolved dependencies}
The comparison displays in Theorem 5.20 and (6.1), p. 24, and again in Section 8, use \(d-k\) where the preceding suspension equivalence requires \(d+k\). At rank one, \(D^1(F)=S^1\) and \(H^{E_1}_{1,0}(\mathbf R_k)=k\), whereas the printed degree \(-1\) right side would be zero. Theorem 6.5’s actual proof already uses \(d+2\). This is recorded as a grading correction rather than a new independent failure of the vanishing theorem. The empty binomial-coefficient ideal at \(n=1\) is \((0)\), consistent with the correctly stated boundary group.\par
Lemma 7.32, p. 48, must refer to coinvariants in its generation claim. With \(A=\mathbb Q\) and rank two, fixing the first apartment \(a\) only produces elements in \(\mathbb Za\otimes\operatorname{St}(\mathbb Q^2)\), not all of \(\operatorname{St}(\mathbb Q^2)\otimes\operatorname{St}(\mathbb Q^2)\). For instance, a distinct apartment \(b\) with a nonzero coefficient on a third line gives \(b\otimes a\) outside that subgroup. The surrounding section and the subsequent applications expressly take coinvariants, where the stated reduction applies. The correction does not refute Corollary 7.33 or Proposition 7.34 in their intended scope.\par
Two items remain Needs Evidence. First, the Moore-complex maps used on p. 61 require torsion representatives, whereas the statement of Proposition 9.3 gives surjectivity onto \(\mathbb Z/2\) and \(\mathbb Z/3\). Surjectivity alone, as illustrated abstractly by \(\mathbb Z\to\mathbb Z/p\), does not supply such representatives. That abstract example is not a counterexample to the actual source construction; the operation construction may supply the missing lift, and no failure is confirmed. Second, Theorem 7.7 is phrased for any connected semi-local ring, while Section 7.5’s resolution argument assumes infinite residue fields. No finite-residue-ring counterexample was established. A section-wide implicit hypothesis is a possible reading, which must be made explicit before asserting an unrestricted extension.\par
The source’s Definition 7.9 says cokernel. A preliminary helper reading of that word as kernel was withdrawn after independent extraction and actual enlarged-page inspection. No version discrepancy or finding is based on that misreading.\par
\section*{Overall impact and state}
Problems 1–3 concern actual derivations on the route to the principal results. Their classification does not assert that all the resulting theorems are false, and the explicit repairs do not lower their core-dependency assessment. The field tensor-square and divided-power arguments survive the local compatible-piece counterexample. Problem 4 gives a false quantitative rank-filtration application even though its finiteness consequence remains valid. The remaining confirmed findings have the narrower scopes stated above.\par
Item | Assessment | Repair verification | Author adoption\par
1 | Major: local product | Changed local model proved in Appendix A | Unknown\par
2 | Major: \(3\)-local proof | Coherent degree calculation checked | Unknown\par
3 | Major: group comparison | Correct extension proof checked | Unknown\par
4 | Major: \(K\) quotient | Kernel formula checked | Unknown\par
5 | Minor: pre-Bloch domain | Admissible set identified | Unknown\par
6 | Minor: finite-field scope | Infinite-field proof retained | Unknown\par
7 | Minor: rank-one boundary | Restricted argument checked | Unknown\par
U1–U2 | Needs Evidence | Not certified | Unknown\par
Scientific decision and adjudication remain pending. No human review or independent acceptance is represented by this report.\par
9 S. Galatius, A. Kupers and O. Randal-Williams, \(E_\infty\)-cells and general linear groups of infinite fields, arXiv:2005.05620v3, acquired 82-page PDF and accompanying source, https://arxiv.org/abs/2005.05620v3. S. Galatius, A. Kupers and O. Randal-Williams, same title, Cambridge accepted manuscript, dated 17 January 2025, 101 pages, repository DOI https://doi.org/10.17863/CAM.115157. H. Matsumoto, Sur les sous-groupes arithmétiques des groupes semi-simples déployés, Ann. Sci. École Norm. Sup. (4) 2 (1969), 1–62, especially pp. 28, 30–31, https://doi.org/10.24033/asens.1174.\par
\section*{Replacement of the restricted local pieces}
Retain the original hypotheses: \(A\) is connected and semi-local with infinite residue fields, and \(M\) is a finitely generated projective module of positive rank. For \(\sigma<E(M)\) define \(W^k_\sigma\) to consist of the basepoint and full lattices \(\phi\) for which \[\phi(a)\in\operatorname{Sub}(M)_\sigma\quad\text{for every }a.\] This is weaker than Definition 7.15. A surviving face or degeneracy only selects or repeats lattice entries, so these are multisimplicial subsets, natural under inclusion of \(\sigma\).\par
Fix a full lattice \(\phi\) and put \(V_a=\phi(a)^\perp\subset M^\vee\). These are finitely generated projective direct summands. The weak compatibility of a simplex \(\tau\) is equivalent to \[\operatorname{span}(\tau\cap V_a)=V_a\quad\text{for every }a.\] Indeed, the clean-evaluation condition makes the span of every subset of \(\tau\) a summand, so taking a double annihilator is legitimate. An adapted dual basis for a splitting of \(\phi\) gives at least one weakly compatible simplex.\par
Consider a finite nonempty collection of weakly compatible simplices \(\tau\). We construct a finite set \(S\) such that every \(\tau\cup S\) remains a simplex and \(S\) spans every \(V_a\) by its elements lying there. For each nonzero \(V_a\), choose \(\operatorname{rank}V_a\) vectors in it successively. For every face \(\rho\) of every current clean simplex \(\tau\cup S\), Lemma 7.11 applies with \(V=V_a\): the containing simplex \(\tau\cup S\) spans \(V_a\) because \(\tau\) already does. Thus only finitely many proper linear subspaces are forbidden modulo each maximal ideal. Also forbid the span of the previously chosen vectors for this \(V_a\), and zero. The residue fields are infinite and there are finitely many maximal ideals, so finite avoidance and the Chinese remainder theorem produce the next vector simultaneously. It is a unimodular functional and adjoining it preserves all required clean intersections. The chosen vectors form a basis modulo every maximal ideal, hence generate \(V_a\) by Nakayama’s lemma. The zero summands need no vectors.\par
At the end, \(\sigma=S\) and every \(\tau\cup\sigma\) are weakly compatible simplices. The inclusions \(\tau\subset\tau\cup\sigma\supset\sigma\) contract every finite subposet of weakly compatible simplices. Its realization is therefore weakly contractible. The fiber over a non-basepoint lattice in the levelwise homotopy colimit is precisely this realization; the basepoint fiber is controlled by the contractibility of \(E(M)\). Consequently \[\operatorname*{hocolim}_{\sigma<E(M)}W^k_\sigma
                      \simeq D^k(M).\]\par
For \(k=1\), \(W^1_\sigma=D^1(M)^\sigma\) with the source’s definition. For \(k=2\), the intersections of two such flags are kernels of unions of subsets of \(\sigma\). Lemma 7.10 makes the relevant sums summands, so the intersection diagram is a full lattice, and every entry belongs to \(\operatorname{Sub}(M)_\sigma\). Hence \[W^2_\sigma\cong
                D^1(M)^\sigma\wedge D^1(M)^\sigma.\] The original one-dimensional connectivity argument now gives \((2d-1)\)-connectivity of each \(W^2_\sigma\) and of their homotopy colimit. Replacing the local pieces by \(W^2_\sigma\) in the quotient comparison on p. 44 likewise restores the displayed resolution equivalence. This repair supplies the missing local-model argument; it does not claim that the strong-compatible pieces themselves have the asserted product decomposition.\par
\section*{Supporting checks}
The source manifest and current-byte inventory accompany this report. Primary PDF SHA-256:\par
4d3a19cc78b188297523396dd8f1e21bdf71986ba6d2894db24ae55c3c05ea0e\par
Accepted PDF SHA-256:\par
fa58cb58cfced3fdfbacb24d37360eda762334f87ef6f40b5c2e3100aab77fc9\par
A whitespace-token comparison of the two extractions is a navigation aid only; it is not a proof of source equivalence. The separate version map records corresponding passages. The primary source extraction was read to zero remaining. Author source files were not compiled or executed.\par
The primary analysis independently reconstructed the rational compatibility counterexample and integrated the helper counterevidence. The exact subset enumeration supports the analytic argument; it is not a general connectivity proof. The broader finite diagnostics for block relations, signs, binomial gcds and the modular centralizer test only their declared finite cases.\par
The current TeX, compiler binary, dependency list, cache inventory, logs, resulting PDF and rendered images are bound by actual hashes in the build and delivery records. Rendering receipts retain their original uninspected flags; actual image observations are separate. A failed Unicode console print during version mapping is preserved; it did not alter either source or the successfully written UTF-8 map.\par
The primary executor and native bounded helpers ran in this conversation. The helpers did not serve as human reviewers, independent acceptors, or replacement controllers.\par
\end{document}
